\section{Discussion}
\label{sec:discussion}

The best one-factor approximation to two correlated traits has a systematic, design-dependent interpretation. With equal item blocks, ML gives an equal composite above a known correlation threshold and favors one trait below it; when the blocks differ, the global solution can favor one trait throughout the correlation range, and an inferior local minimum can favor the other. With unequal block lengths, the favored trait can also switch abruptly at a correlation where two different solutions fit equally well. These changes in representation are accompanied by exact changes in population fit, so the formulas make it possible to examine both what the fitted factor represents and what a fit criterion can distinguish for a specified measurement design.

\subsection{Interpreting One-Factor Fit}

The cutoff boundaries in Table~\ref{tab:1} quantify the limits of an inference from acceptable fit to unidimensionality: each gives the trait correlation from which a conventional cutoff is met although two traits generate the data. Every CFI boundary in Table~\ref{tab:1} lies above the ML transition, so in these equal-block designs the approximation meeting the CFI cutoff is a composite. This need not hold for every index: with loading .50 and four items, the SRMR cutoff is met at every correlation, including the tilted regime. A tilted solution and an acceptable fit are therefore separate findings, although unequal blocks can make trait representation unbalanced even at high correlations. Neither the fit index nor the number of fitted factors alone specifies the contribution of each trait, and a fit index is not a measure of the substantive importance of a trait.

Rather than a universal replacement cutoff, the analysis supplies a benchmark conditional on the proposed item structure, which can be reported alongside the observed indices and examined over plausible loading values. Such a calculation makes the evidential limits of a merged-factor comparison explicit, particularly when absolute fit of the merged model is treated as evidence against discriminant validity (R\"onkk\"o \& Cho, 2022). It complements simulation-based dynamic fit indices, whose one-factor implementation considers residual correlations as the source of misspecification (McNeish \& Wolf, 2023), and the simulation-cum-ROC approach of Groskurth et al.\ (2025), which permits a two-factor alternative to a one-factor model and selects cutoffs from simulated classification performance. The boundaries derived here identify the population trait correlations at which a fixed cutoff is met; they do not calibrate finite-sample type I and type II error rates.

Interpretation also requires the loading pattern, the intended score and the estimator. For a specified two-trait generating model, the fitted loadings determine the trait contributions only together with the generating loadings; raw one-factor loadings alone do not identify the factor's trait composition. The composition of an estimated factor must also be distinguished from that of a score computed from the items: regression and Bartlett factor scores have the same relative trait contributions in the setting considered here, but these generally differ from those of the fitted factor and the unit-weighted sum. Reporting one score therefore requires an interpretation of that score's weights as well as an assessment of model fit. ML, unweighted least squares, and principal components can represent different directions for the same population correlation matrix, so the estimation criterion matters beyond standard errors and goodness-of-fit corrections. Likewise, the close agreement of lavaan's robust categorical indices (Savalei, 2021) with the continuous-ML references near the cutoffs is a numerical result for the studied designs and does not identify the categorical fitted factor with the ML fitted factor.

\subsection{Computation and Reporting}

In the finite-sample study, item order, an arbitrary feature of the input, was associated with which trait the default-start solution favored and, when the blocks differed, with whether that solution had a higher objective and favored the weaker block, consistent with lavaan's default starting values, which are anchored on the first-listed item (Section~\ref{sec:numerical}). For a one-factor model fitted to known item blocks, a useful check is therefore to compare the solutions obtained with each block listed first and with additional starting values, retaining the discrepancy and loading pattern of each distinct solution. The archived multistart results support this check for the studied designs (Section~\ref{sec:numerical}) but do not prove that two orders are sufficient in general. Moreover, agreement in objective and reported convergence do not by themselves establish stationarity or a local minimum: the population refits showed that equal starting covariances could lead to different trajectories under different identifications, including a symmetric saddle. When competing solutions affect interpretation, objective comparisons should be accompanied by checks of the relevant first- and second-order conditions.

The boundary criterion of Proposition~\ref{prop:boundary} (Appendix~\ref{app:heywood-proof}) has a separate use. For a one-factor ML model with nonnegative uniquenesses, it gives a necessary condition for a solution with one zero uniqueness to be a local minimum. In the two-trait population it excludes such minima except at the stated two-item, zero-correlation endpoint; more generally, it can rule out a reported zero-uniqueness solution when releasing that uniqueness decreases the discrepancy. The result does not classify unconstrained negative uniqueness estimates or solutions obtained by other estimators.

The accompanying R package \texttt{undeRfactor} (version 0.2.0, \url{https://github.com/atomashevic/undeRfactor}) implements the population calculations, the ML, ULS and WLSMV transition thresholds (\texttt{phi\_star()}), the ML cutoff boundaries (\texttt{blind\_spot()}), all stationary points of designs with one loading per block from the quartic of Appendix~\ref{app:unequal-scalar} (\texttt{block\_solutions()}), the candidate-solution search, the coefficient and score shares, the one-factor solution implied by an estimated two-factor model (\texttt{predict\_one\_factor()}) and the item-order check. Closed forms and the quartic are used where available; elsewhere, the package reports the distinct candidates found and does not certify global optimality. Coefficient shares describe covariance matching rather than explained variance; signed shares depend on trait orientation and are undefined when the coefficients sum to zero, and nonzero reference block means and $|\varphi|<1$ are required for identification.

\subsection{Scope and Further Work}

The exact results concern a structured two-trait population. The exchangeable case permits a closed global solution, and constant loadings within unequal blocks permit an exact reduction to a quartic equation. For unequal block lengths, this reduction gives a selection rule but not a closed partition of the full design space, and the global solution need not vary continuously with the trait correlation (Section~\ref{sec:unequal}). For loadings that vary within blocks, the general information-loss bound (Appendix~\ref{app:ceiling}) and numerical examples provide partial results. Independent residuals are part of the generating model, and more than two traits are outside the present analysis. These restrictions isolate a mechanism that is difficult to identify from numerical optimization alone, but they also limit direct application of the formulas to heterogeneous instruments. Matching lower bounds for variable loadings, a full classification of competing minima with unequal block lengths, and extensions to more than two traits are therefore priorities for further analytic work.

The population and finite-sample conclusions have different scopes: the normal-theory ML theorems concern the population objective, whereas the simulation study covers sample sizes from 250 to 2,000 and the specified continuous and ordinal designs. The study does not establish performance for smaller samples, nonnormal continuous responses, or every set of ordinal thresholds; in particular, different item thresholds generally remove the two-weight structure used in the weighted least-squares result. Archived solution groups identify alternatives found by the specified starts but do not certify sample stationarity or global optimality. The power summaries use conventional noncentral-$\chi^2$ approximations; their finite-sample accuracy, especially near the transition and at nonregular endpoints, was not assessed.

The cross-loading and application results require corresponding qualifications. Equal cross-loadings yield an exact transformation of the latent-response covariance into the two-block form. The Gaussian analogue of the Strachan et al.\ (2022) design suggests a mechanism for their reported instability at negative correlations but is not a replication of binary marginal-ML estimation, and the binary calculations (Appendix~\ref{app:binary}) provide numerical evidence rather than a general theorem for IRT models. Likewise, the PHQ-ADS illustration applies normal-theory ML to the published two-factor loadings, which were estimated from ordered categorical responses with a logistic link, and uses a common-loading range of .70 to .80 as a sensitivity analysis. It is therefore not an exact reanalysis of the original categorical responses and cannot by itself validate the total score.

The exchangeable two-trait population also admits a bifactor representation with general loadings $\lambda\sqrt{\varphi}$ and specific loadings $\lambda\sqrt{1-\varphi}$, whose explained common variance is therefore $\varphi$, linking the present thresholds to an index used to assess essential unidimensionality (Reise et al., 2013; Rodriguez et al., 2016). The present results already show why one-factor fit should be reported together with the design, the estimator, and the trait representation of the solution that was obtained.

% The named preprint carries the full Declaration of AI Use (declarations.tex).
\ifdefined\preprint\else
\paragraph{Use of AI tools.}
The mathematical proofs reported in this article and in Online Resources 1--3 were developed with the assistance of the large language model ChatGPT 6 Pro (OpenAI). The large language model Claude Opus 5.5 (Anthropic) assisted with writing the code and with preparing and editing the manuscript.
\fi
